% © 2026 Ing. David Jaroš — CC BY-NC-ND 4.0
% UBT-AI-PROVENANCE-BEGIN
% schema: ubt-ai-provenance/v1
% tier: C_working
% ai_assistance: disclosed
% human_review: risk-based
% editorial_responsibility: Ing. David Jaroš
% policy: ../../AI_PROVENANCE.md
% notice: Working material; exhaustive human review is not claimed.
% UBT-AI-PROVENANCE-END

\documentclass[11pt,a4paper]{article}
\usepackage[margin=2.5cm]{geometry}
\usepackage{amsmath,amssymb,amsthm,mathtools}
\usepackage[most]{tcolorbox}
\newtheorem{theorem}{Theorem}
\newtheorem{corollary}[theorem]{Corollary}
\newtheorem{remark}[theorem]{Remark}
\title{No-Extra-Variable Rank Theorem for the Canonical UBT Generalized-Dirac Relation}
\author{Ing. David Jaro\v{s}}
\date{27 July 2026}
\begin{document}
\maketitle

\begin{abstract}
The canonical UBT metric is built from the single master field through
$E_\mu=\mathcal N_0^{-1/2}D_\mu\Theta$ and
$g_{\mu\nu}=e_\mu{}^ae_\nu{}^b\eta_{ab}$.  The unconstrained tetrad-to-metric
map has rank ten and a six-dimensional Lorentz kernel.  This note proves the
exact rank formula after imposing four complex generalized-Dirac equations,
without introducing any auxiliary variable or second field.  Full metric rank
is retained exactly when the admissible tetrad tangent space is transverse to
the Lorentz kernel.  In particular, an invertible Jacobian of the equation
with respect to the value of the original field makes all tetrad variations
admissible.  A nonzero scalar or scalar--pseudoscalar zero-order Dirac block
provides an explicit sufficient realization, including under strict complex-
time holomorphy.  Conversely, eight independent real constraints acting only
on the tetrad reduce the metric rank to at most eight.
\end{abstract}

\section{Canonical variables}
The only field is
\[
 \Theta(q,\tau)\in\mathbb C\otimes\mathbb H,
 \qquad \tau=t+i_{\rm c}\psi,
\]
written equivalently as $\Psi=\operatorname{vec}(\Theta)\in\mathbb C^4$.
On the Lorentz slice,
\[
 E_\mu=\mathcal N_0^{-1/2}D_\mu\Theta
 =i_{\rm c}e_\mu{}^0\mathbf1+e_\mu{}^k\mathbf e_k,
\]
so $e=(e_\mu{}^a)\in\mathbb R^{16}$ and
\[
 g=e\eta e^T,
 \qquad \eta=\operatorname{diag}(-1,1,1,1).
\]
At every nondegenerate tetrad,
\[
 \operatorname{rank}D_eg=10,
 \qquad K:=\ker D_eg,\quad\dim K=6.
\]

\section{Exact constrained-rank formula}
Let the real form of the four-complex-component field equation be
\[
 F(\Psi,e)=0\in\mathbb R^8.
\]
At a solution denote its Jacobian blocks by $F_\Psi$ and $F_e$.  A tetrad
variation is compatible with the linearized equation precisely when its
change can be cancelled by a variation of the value of the same original
field.  Thus
\[
 \mathcal A:=\{\delta e:\ F_e\delta e\in\operatorname{im}F_\Psi\}.
\]

\begin{theorem}[Constrained metric rank]
For every nondegenerate tetrad,
\[
 \boxed{
 \operatorname{rank}(D_eg|_{\mathcal A})
 =\dim(\mathcal A+K)-6.}
\]
Consequently the constrained metric rank is ten if and only if
\[
 \boxed{\mathcal A+K=\mathbb R^{16}.}
\]
\end{theorem}

\begin{proof}
Since $D_eg$ is surjective with kernel $K$, it identifies
$\mathbb R^{16}/K$ with the ten-dimensional metric tangent space.  The image
of $\mathcal A$ is therefore isomorphic to $(\mathcal A+K)/K$.  Taking
dimensions proves the formula and the equivalence.
\end{proof}

\begin{corollary}[Original-field absorption]
If $F_\Psi$ is invertible as an $8\times8$ real matrix, then
$\mathcal A=\mathbb R^{16}$ and
\[
 \operatorname{rank}(D_eg|_{F=0})=10.
\]
For every $\delta e$, the unique compatible variation is
\[
 \delta\Psi=-F_\Psi^{-1}F_e\delta e.
\]
\end{corollary}

This uses no added variable: $\delta\Psi$ is the variation of the value of the
same UBT master field whose first derivative supplies the tetrad.

\begin{corollary}[Eight-constraint no-go]
If $F_\Psi=0$ and $F_e$ has real rank eight, then
$\dim\mathcal A=8$ and
\[
 \operatorname{rank}(D_eg|_{\mathcal A})\le8.
\]
Thus four independent complex equations cannot act only on the tetrad and
still retain generic metric rank ten.
\end{corollary}

\begin{remark}
Constraints need not lower the metric rank when they remove only Lorentz-frame
directions.  The exact condition is transversality $\mathcal A+K=\mathbb R^{16}$,
not merely a count of equations.
\end{remark}

\section{Holomorphic generalized-Dirac application}
The canonical Clifford lift gives
\[
 \Gamma_\mu=\mathcal C(E_\mu),
 \qquad \tfrac12\{\Gamma_\mu,\Gamma_\nu\}=g_{\mu\nu}I_4,
\]
and an anticommuting fifth channel $\Gamma_\psi$.  Consider
\[
 i_{\rm c}\Gamma^\mu D_\mu\Psi
 +i_{\rm c}\Gamma_\psi D_\psi\Psi
 -\mathcal M\Psi=0.
\]
Strict covariant holomorphy gives $D_\psi\Psi=i_{\rm c}D_0\Psi$, while
$D_\mu\Psi=\sqrt{\mathcal N_0}\operatorname{vec}(E_\mu)$.  Therefore
\[
 \mathcal M\Psi
 =\sqrt{\mathcal N_0}\left[
 i_{\rm c}\Gamma^\mu(e)\operatorname{vec}(E_\mu)
 -\Gamma_\psi\operatorname{vec}(E_0)
 \right].
\]
For $\mathcal M=mI_4$, $m\ne0$, this equation uniquely determines $\Psi$ for
every admissible tetrad.  Its complex determinant is $m^4$, hence its real
Jacobian has rank eight.  The preceding corollary proves pointwise metric rank
ten on the first-jet equation manifold.

For
\[
 \mathcal M=m_sI_4+i_{\rm c}m_p\Gamma_*,
 \qquad \Gamma_*^2=I_4,
\]
\[
 \mathcal M^{-1}
 =\frac{m_sI_4-i_{\rm c}m_p\Gamma_*}{m_s^2+m_p^2},
 \qquad
 \det_{\mathbb C}\mathcal M=(m_s^2+m_p^2)^2.
\]
The realified determinant is $(m_s^2+m_p^2)^4$, so the same conclusion holds
whenever $(m_s,m_p)\ne(0,0)$.

\begin{tcolorbox}[colback=green!5!white,colframe=green!50!black,title=Exact verdict]
The rank problem has an exact no-extra-variable solution at the pointwise
first-jet level: the finalized equation must satisfy
$\mathcal A+K=\mathbb R^{16}$.  An invertible equation Jacobian with respect
to the original field value is a sufficient condition and a nonzero Dirac
zero-order block realizes it explicitly.  The remaining task is dynamical,
not a rank count: derive that block from the canonical UBT action and prove
local existence/integrability of the resulting implicit holomorphic PDE.
\end{tcolorbox}

\end{document}
